\documentclass[10pt,a4paper]{amsart}
\usepackage[margin=19mm]{geometry}
\usepackage{amsmath,amssymb,mathtools}
\usepackage[scr=rsfs,cal=euler]{mathalfa}
\usepackage{xfrac}
\usepackage[unicode,hidelinks,bookmarks=false]{hyperref}
\newcommand{\ie}{i.e.\ }
\newcommand{\cf}{cf.\ }
\newcommand{\resp}{respectively\ }
\newcommand{\Subsection}{\subsection}
\providecommand{\nobreakdash}{\nobreak\hbox{-}}
\setlength{\emergencystretch}{2em}
\multlinegap0pt
\usepackage{amssymb}
\usepackage{mathtools}
\usepackage[shortlabels]{enumitem}
\newtheorem{prop}{Proposition}
\newtheorem{cor}{Corollary}
\title{Cyclotomic Numbers of Order $q-1$ over $\mathbb{F}_{q^r}$}
\author{Hayaki Kudo, Yuto Nogata}
\date{Source: arXiv:2604.25461v1 (CC BY 4.0)}
\begin{document}
\maketitle

\noindent\emph{Source and licence.} Cyclotomic Numbers of Order $q-1$ over $\mathbb{F}_{q^r}$. Version arXiv:2604.25461v1 (CC BY 4.0), distributed by arXiv under the Creative Commons Attribution 4.0 International licence, http://creativecommons.org/licenses/by/4.0/, as recorded in the arXiv licence metadata for this exact version and shown on the abstract page https://arxiv.org/abs/2604.25461v1. Full licence notice: \texttt{licenses/arxiv-2604.25461-CC-BY-4.0.txt}.\par\smallskip

\noindent\textit{Editorial scope.} Counted unit: the article's cor cor:character-sum-half-k-qgefour-rgeqthree (source0.tex lines 596--598). The article's complete original proof of it is delivered with it (source0.tex lines 600--621, one \texttt{proof} environment of 22 lines). No mathematics is written, repaired or re-derived: the statement and the proof are the article's own lines, in the article's own notation, and no retained line is substituted or re-flowed.  Retained uncounted as the source-local context the proof needs: lines 126--133; lines 451--479.  Nothing was omitted from any retained span, so the package carries no omission record.  No retained line contains a citation, so the wrapper carries no bibliography and no bibliography entry.  No retained line contains a figure environment, an image-inclusion command or drawing code, so no image is delivered and the package declares no asset dependency.  Editorial formatting only, in addition: an AMS \texttt{amsart} preamble that restates the article's own declarations and macros used by the retained lines, namely the environments \texttt{prop}, \texttt{cor} on separate counters, together with the standard packages those lines need.  The retained environments print in this document as Proposition 1, Corollary 1 in order of appearance, whereas the article prints the same items with the numbering of its own full document.  The complete native source of the article as distributed in the arXiv e-print for this exact version is bundled unchanged as \texttt{sources/199305/source0.tex}.
\begin{equation}\label{eq:ceil-half-k}
\left\lceil \frac{k}{2}\right\rceil=
\begin{cases}
\frac{q^r+q-2}{2(q-1)} & \text{if $q$ is even,} \\
\frac{q^r+q-2}{2(q-1)} & \text{if $q$ is odd and $r$ is odd,} \\
\frac{q^r-1}{2(q-1)} & \text{if $q$ is odd and $r$ is even.}
\end{cases}
\end{equation}

\begin{prop}\label{prop:main-term-explicit-general}
For $0\le a,b\le q-2$, define
\[
  A_r(a):=
  \begin{cases}
    -(q-2), & \textup{if } r \textup{ is even and } a=0,\\
    -(q-2), & \textup{if } r \textup{ is odd},\ q \textup{ is even, and } a=0,\\
    -(q-2), & \textup{if } r \textup{ is odd},\ q \textup{ is odd, and } a=\frac{q-1}{2},\\
    1, & \textup{otherwise,}
  \end{cases}\]
  \[
  B(b):=
  \begin{cases}
    -(q-2), & \textup{if } b=0,\\
    1, & \textup{if } b\neq 0,
  \end{cases}\quad
  C(a,b):=
  \begin{cases}
    -(q-2), & \textup{if } a=b,\\
    1, & \textup{if } a\neq b.
  \end{cases}
\]
Then
$(a,b)_{q-1}=M_{a,b}+E_{a,b}$,
where
$M_{a,b}:=\frac{q^r-2+A_r(a)+B(b)+C(a,b)}{(q-1)^2}$
and
$|E_{a,b}|\le \frac{(q-2)(q-3)}{(q-1)^2}q^{r/2}$.
\end{prop}

\begin{cor}\label{cor:character-sum-half-k-qgefour-rgeqthree}
Assume that $q\ge4$ and $r\ge3$. Then $(a,b)_{q-1}< \left\lceil \frac{k}{2}\right\rceil$ for every $0\le a,b\le q-2$.
\end{cor}

\begin{proof}
By Proposition \ref{prop:main-term-explicit-general}, we have
$(a,b)_{q-1}\le \frac{q^r+1}{(q-1)^2}+\frac{(q-2)(q-3)}{(q-1)^2}q^{r/2}$.

\noindent(i): Assume that $q$ is even or $r$ is odd. Then \eqref{eq:ceil-half-k} gives $\left\lceil \frac{k}{2}\right\rceil=\frac{q^r+q-2}{2(q-1)}$, and hence
\[
\left\lceil \frac{k}{2}\right\rceil-(a,b)_{q-1}
\ge
\frac{q-3}{2(q-1)^2}\left(q^r-2(q-2)q^{r/2}+q\right).
\]
Since $r\ge3$ and $q\ge4$, we have $q^{r/2}\ge q^{3/2}>2(q-2)$. Thus $q^r-2(q-2)q^{r/2}+q>0$, and the desired result follows.

\noindent(ii): Assume that $q$ is odd and $r$ is even. Then $q\ge5$ and $r\ge4$. By \eqref{eq:ceil-half-k}, we have $\left\lceil \frac{k}{2}\right\rceil=\frac{q^r-1}{2(q-1)}$, and hence
\[
\left\lceil \frac{k}{2}\right\rceil-(a,b)_{q-1}
\ge
\frac{q^{r/2}(q-3)\left(q^{r/2}-2(q-2)\right)-(q+1)}{2(q-1)^2}.
\]
Since $q^{r/2}\ge q^2$ and $q^{r/2}-2(q-2)\ge q^2-2q+4$, we obtain
$(q-3)\left(q^r-2(q-2)q^{r/2}\right)\ge (q-3)q^2(q^2-2q+4)>q+1$.
Therefore $(a,b)_{q-1}< \left\lceil \frac{k}{2}\right\rceil$.
\end{proof}

\end{document}
